% !TeX root = ../../packet-and-principal-models.tex
\section{Integral models, polarizations and maps}
We first give coordinates for the common principal component. We then
separate marked, polarized and unpolarized equivalence. The ideal and unit
calculations are also the explicitly cited arithmetic inputs to Paper II.


\begin{proposition}[An explicit principal RM component]
\label{pp:prop:explicit}
Write $O=\OL$, $u=t-1$, $d=f'(t)$, and
\[
 \beta=-t^5+t^4+6t^3-2t^2-8t-2,\qquad
 a=(\beta u^2)^{-1},\quad b=u^{-3},\quad
 \delta=ab=(\beta u^5)^{-1}.
\]
The common self-dual lattice $P=\Lambda_{\mathrm{common}}$ has $O$-basis
\begin{equation}\label{pp:eq:basis}
 T=\begin{pmatrix}a&3b\\0&b\end{pmatrix},\qquad P=TO^2.
\end{equation}
Its determinant $\delta$ generates the codifferent and has signs
$(+,-,+,-,+,-)$ in the ordered embeddings $1,5,11,13,17,19$.
A principal RM presentation of the family in
Theorem~\ref{thm:principalfamily} has homology module
\[
 P_{\rm std}=Oe_1+\delta Oe_2
\]
with the trace form $E_1$, period lattice $Ow+\delta O$, and arithmetic group
\begin{equation}\label{pp:eq:group}
 \Gamma_{\rm std}
 =\left\{\begin{pmatrix}r&\delta s\\v/\delta&q\end{pmatrix}:
 r,s,v,q\in O,\ rq-sv=1\right\}
 =\operatorname{diag}(\delta,1)\SL_2(O)
   \operatorname{diag}(\delta,1)^{-1}.
\end{equation}
Its period map on the universal covering curve is
\begin{equation}\label{pp:eq:curve}
 F_{{\rm std},i}(z)
 =\frac{\sigma_i(a^2)F_{t,i}(z)}{3F_{t,i}(z)+1}.
\end{equation}
Thus its image lies in the specified principal RM component
$\HH^6/\Gamma_{\rm std}$, with a compatible neat refinement when a fine
moduli interpretation is required.
\end{proposition}
\begin{proof}
The ideal identities $(u)=\mathfrak p_7$, $(\beta)=\mathfrak p_3$ and
$(d)=\mathfrak p_3\mathfrak p_7^5$ give the unit
\[
 \varepsilon=\frac{d}{\beta u^5}
 =28+215t+335t^2+82t^3-87t^4-35t^5,
 \qquad N_{L/\Ql}(\varepsilon)=1.
\]
They may also be checked directly from $f$: $N(u)=7$,
$N(\beta)=-27$ and $\beta$ vanishes modulo the degree-three factor at $3$.
Under multiplication by $d$, the local conditions defining the common
subgroup in Proposition~\ref{prop:commonprincipal} say
\[
 y\in\beta O,\qquad y\in u^2O,\qquad x-3y\in u^3O.
\]
Since $(\beta)$ and $(u)$ are coprime, their simultaneous solution is
\[
 dP=Ou^3(1,0)+O\beta u^2(3,1).
\]
Now
\[
 dT=\varepsilon
 \begin{pmatrix}u^3&3\beta u^2\\0&\beta u^2\end{pmatrix},
 \qquad
 T^{-1}=\begin{pmatrix}\beta u^2&-3\beta u^2\\0&u^3\end{pmatrix}.
\]
The unit factor leaves the module unchanged, proving
\eqref{pp:eq:basis} and $O^2\subset P$. Also $\delta d=\varepsilon$,
so $\delta O=\mathfrak D_L^{-1}$ and
$[P:O^2]=|N(\delta)|^{-1}=453789$.

In the basis $T$, the pairing is
$E_\delta((r,s),(r',s'))=\operatorname{Tr}_{L/\Ql}
(\delta(r's-rs'))$. Exact arithmetic gives
\[
 \delta=-\frac{11}{21}-\frac{209}{21}t-\frac{55}{3}t^2
       -\frac{110}{21}t^3+\frac{33}{7}t^4+2t^5.
\]
The integral matrix $(\operatorname{Tr}(\delta t^{i+j}))_{0\leq i,j\leq5}$
has determinant $-1$, so its alternating trace form is unimodular.
Disjoint rational isolating intervals for all six roots of $f$, followed
by rational interval evaluation of this polynomial, certify the stated
signs. The exact identities and signs are checked in
\path{experiments/principal_normalization_v111/explicit_basis.py}
and separately, using rational polynomial arithmetic, in
\path{experiments/principal_normalization_v111/review/independent_basis_check.py}.

At a period $z\in\HH^6$, interpreted componentwise, the basis columns
evaluate to $az$ and $b(3z+1)$. Division by $b(3z+1)$ gives the period
\[
 \zeta=\frac{a}{b}\frac{z}{3z+1}
\]
on the standard $O^2$ module with form $E_\delta$. Its imaginary-part
signs are those of $\delta$; $E_\delta$ is positive on this mixed
component. Multiplication of the complex coordinate by $\delta$ gives
$w=\delta\zeta=a^2z/(3z+1)$ and lattice $Ow+\delta O$. The corresponding
homology identification is
\[
 C=\operatorname{diag}(1,\delta)T^{-1},\qquad \det C=1.
\]
It maps $P$ onto $P_{\rm std}$ and preserves $E_1$ exactly. Thus this
is a polarized isomorphism, not just an isogeny comparison.
The principal module $P_{\rm std}$ is self-dual because $\delta O$ is
the codifferent. Its transition stabilizer for $g\mapsto g^{-\mathsf T}$
is precisely \eqref{pp:eq:group}. The form $E_1$ is positive on its
all-upper-half-plane domain.

The composite period matrix is
\[
 W=\begin{pmatrix}a^2&0\\3&1\end{pmatrix}\in\GL_2^+(L),
 \qquad \det W=a^2.
\]
It gives \eqref{pp:eq:curve}; its denominators do not vanish on $\HH$.
The stabilizer calculation below verifies equivariance on the same
connected cover as Theorem~\ref{thm:principalfamily}. That theorem's
algebraization identifies the resulting moduli map with the original
principal family, preserving its degree-$453789$ isogeny and its three
specialized $K$-actions.
\end{proof}

\begin{proposition}[Explicit transition stabilizer]\label{pp:prop:stabilizer}
For $g=\left(\begin{smallmatrix}r&s\\v&q\end{smallmatrix}\right)
\in\SL_2(O)$ put $Q(g)=v+3r-3q-9s$. Then
\[
 \Gamma_H=\{g\in\SL_2(O):s\in\beta O,\ Q(g)\in uO\},
 \qquad
 W\Gamma_HW^{-1}=\Gamma_{\rm std}\cap W\SL_2(O)W^{-1}.
\]
\end{proposition}
\begin{proof}
The homology basis matrix is $T$, the period matrix is $S=T^{\mathsf T}$,
and the transition action is $g^{-\mathsf T}$. Direct multiplication gives
\[
 T^{-1}g^{-\mathsf T}T
 =\begin{pmatrix}
 q+3s&-(\beta/u)Q(g)\\-(u/\beta)s&r-3s
 \end{pmatrix}.
\]
Its integrality is equivalent to the two displayed congruences, because
$u$ and $\beta$ are coprime. Determinant one then gives a lattice
automorphism. On the final periods,
\[
 WgW^{-1}
 =\begin{pmatrix}r-3s&a^2s\\a^{-2}Q(g)&q+3s\end{pmatrix}.
\]
The identities $a^2/\delta=u/\beta$ and $\delta a^{-2}=\beta/u$ prove
the intersection assertion. Consequently \eqref{pp:eq:curve} is
equivariant for $Wj(\Delta')W^{-1}\subset\Gamma_{\rm std}$.
The indices $224$ and $28$ are those of
Proposition~\ref{prop:principalstabilizer}; no new degree after neat
refinement is inferred.
\end{proof}

\subsection{Marked and unmarked distinctions}
\begin{proposition}[Distinct principal points with fixed real multiplication]
\label{pp:prop:markeddistinct}
Let $B_X$, for $X=A,B,C$, be the three specified principal fibers of
Theorem~\ref{thm:principalfamily}, with their fixed $\OL$ actions.
For $X\ne Y$ one has
\[
 \operatorname{Hom}_{\OL}(B_X,B_Y)\otimes_{\mathbf Z}\Ql=0.
\]
Consequently their images in the principal RM component
$\HH^6/\Gamma_{\rm std}$ are pairwise distinct, also after forgetting
auxiliary level while retaining the $\OL$ marking.
\end{proposition}
\begin{proof}
The actual geometric CM types, expressed using $z=\zeta_{42}$, are
\[
 \Phi_A=\{1,5,17,19,29,31\},\quad
 \Phi_B=\{1,5,23,25,29,31\},\quad
 \Phi_C=\{1,5,11,13,17,19\}.
\]
Their $M$ actions restrict to the same prescribed $L$ action:
the generators for $B,C$ satisfy $Z^2-tZ+I=0$, and the certified
generator $Z_A=(A_t+tI)/(t^2-1)$ of
Proposition~\externalref{supplement}{prop:Amaximalpair} satisfies the same identity and
$Z_A^3=A_t$. Passing to transposes gives the homology actions without
changing the scalar $t$. Thus $z+z^{-1}$ acts as the fixed scalar $t$
in every case. These are the normalized geometric types, not companion types
or representatives re-marked by arbitrary Galois automorphisms.
The lattice inclusions of Theorem~\ref{thm:principalfamily} are the
identity on rational $L^2$, and the coordinate normalization in
Proposition~\ref{pp:prop:explicit} commutes with $L$ at each real
place. Hence the same marked rational Hodge structures, with their
rational $M$ actions, describe the principal fibers.

For completeness, consider two rank-one $M$ rational Hodge structures
with this fixed $L$ restriction. After identifying each rational
vector space with $M$, every rational linear map is uniquely
\[
 F=\sum_{g\in G}m_g g,\qquad m_g\in M,\qquad
 G=\operatorname{Gal}(M/\Ql).
\]
Indeed, after splitting $M$ over $\mathbf C$, each summand $Mg$
occupies the matrix positions of its embedding permutation. These
supports are disjoint and their dimensions sum to $[M:\Ql]^2$.
Commutation with every $a\in L$ gives
$m_g(g(a)-a)=0$. Since $M$ is a field, only
$g\in\operatorname{Gal}(M/L)=\{1,c\}$ can remain, where $c$ is
complex conjugation. If $m_g\ne0$, every conjugate of $m_g$ is
nonzero. Its entries cannot cancel those of another summand.
Preserving the Hodge subspaces therefore requires the two types to
be equal for $g=1$, or complementary for $g=c$.

The three displayed types are pairwise unequal, and all contain $1$
and omit $41$, so no two are complementary. Thus every $L$-linear
Hodge map between distinct fibers is zero. An $\OL$-linear
homomorphism of abelian varieties induces such a map on rational
homology, and a homomorphism with zero map on homology is zero.
This proves the assertion without a primitivity assumption.
An identification in the specified principal RM moduli component
would be an $\OL$-linear polarized isomorphism, which is excluded.
The unmarked isogeny relation $19\Phi_B=\Phi_C$ from
Proposition~\ref{prop:knownCM} changes the $L$ marking and gives no
contradiction. The next proposition treats the unmarked principal
images; neither result identifies Mostaed's named image.
\end{proof}


\begin{proposition}[Three distinct unmarked principal fibers]
\label{pp:prop:unmarkeddistinct}
The three specified principally polarized fibers $B_A,B_B,B_C$ of
Theorem~\ref{thm:principalfamily} are pairwise nonisomorphic as principally
polarized complex abelian varieties, after forgetting both the fixed
$\OL$ action and auxiliary level. Thus they give three distinct points
of the coarse moduli space of principally polarized abelian sixfolds.
\end{proposition}
\begin{proof}
By Proposition~\ref{prop:knownCM}, the underlying $A$ fiber is isogenous
to $J(y^2=x^7-1)^2$, whereas the $B,C$ fibers are isogenous to the simple
sixfold $X_{21}$. The principal-model isogenies preserve this distinction,
so only $B_B$ versus $B_C$ remains.

Write $z=\zeta_{42}$, $t=z+z^{-1}$, $u=t-1$ and
\[
 \beta=-t^5+t^4+6t^3-2t^2-8t-2,\qquad
 a=(\beta u^2)^{-1},\quad b=u^{-3}.
\]
The actual common lattice of Proposition~\ref{pp:prop:explicit} is
$P=\OL(a,0)+\OL(3b,b)$, with the restriction of $E_1$.
For $X=B,C$, let $(c_X,d_X)$ be the bottom row of the normalized
vertex matrix $X_t$. Explicitly,
\[
 c_B=t^{-1},\quad d_B=t,\qquad
 c_C=t^{-1}(-1+5t+4t^2-t^3-t^4),\quad
 d_C=-t+4t^2+4t^3-t^4-t^5.
\]
On homology, $z$ acts by $X_t^T$. The identification
$J_X:M\longrightarrow L^2$ given by
$J_X(v_0+v_1z)=(c_Xv_1,v_0+d_Xv_1)$ therefore has
\[
 I_X=J_X^{-1}(P)=\OL a w_X+\OL b(1+3w_X),\qquad
 w_X=(z-d_X)/c_X.
\]
In these coordinates its actual polarization is
\[
 E_X(v,w)=\Tr_{M/\Ql}(\xi_Xv\bar w),\qquad
 \xi_X=-\frac{c_X}{z-z^{-1}}.
\]
Indeed the relative trace equals $c_X(v_0w_1-v_1w_0)$, which
recovers $E_1(J_Xv,J_Xw)$ after applying $\Tr_{L/\Ql}$.

The disjoint-support argument in
Proposition~\ref{pp:prop:markeddistinct} applies without imposing
$L$-linearity: a term $m_g\sigma_g$ can preserve the Hodge subspaces
only if $\Phi_B=g\Phi_C$. For the actual types displayed there, the
unique such exponent is $g=31$ modulo $42$. Hence every possible
isomorphism has the form $F(v)=m\sigma_{31}(v)$. Its integrality and
polarization conditions imply, respectively,
\[
 m\in J:=\frac{I_C}{\sigma_{31}(aw_B)},\qquad
 m\bar m=r:=\frac{\sigma_{31}(\xi_B)}{\xi_C}.
\]
The first condition uses only the image of the first actual lattice
generator; it requires no invertibility of a nonmaximal-order ideal.
The second follows by moving $\sigma_{31}$ through the absolute trace.
In particular $q(m):=\Tr_{M/\Ql}(m\bar m/r)=12$ is necessary.

We record the finite obstruction explicitly. Let $Q$ be the rational
Gram matrix of $q$ on the basis of $J$ induced by
$t^ja w_C,t^jb(1+3w_C)$, $0\leq j<6$, divided by
$\sigma_{31}(aw_B)$. The certificate
\path{certificates/principal_unmarked_bc_v122.py} reconstructs this
matrix from the displayed field and polarization formulas. Its data
file \path{certificates/principal_unmarked_bc_data.json} supplies an
integral matrix $U$ with determinant $-1$ and the exact identity
$U^TQU=H$, where
\begingroup
\setlength{\arraycolsep}{2.3pt}
\scriptsize
\[
H=\left(\begin{array}{*{12}{r}}
14&-7&0&0&-6&-2&3&2&0&0&-3&6\\
-7&14&0&0&3&4&-6&-1&0&0&-3&-3\\
0&0&14&0&4&-1&-2&-6&0&-7&2&-4\\
0&0&0&14&-1&-6&2&-2&-7&0&1&1\\
-6&3&4&-1&16&4&-8&-8&2&-2&-1&2\\
-2&4&-1&-6&4&16&-8&4&3&2&2&2\\
3&-6&-2&2&-8&-8&16&4&-1&4&-1&-1\\
2&-1&-6&-2&-8&4&4&16&4&3&-1&2\\
0&0&0&-7&2&3&-1&4&14&0&1&-2\\
0&0&-7&0&-2&2&4&3&0&14&2&2\\
-3&-3&2&1&-1&2&-1&-1&1&2&22&-11\\
6&-3&-4&1&2&2&-1&2&-2&2&-11&22
\end{array}\right).
\]
\endgroup
Exact rational factorization proves $H$ positive definite, and direct
inversion gives $12(H^{-1})_{ii}<4$ for every $i$. Cauchy--Schwarz
therefore places every integer vector of norm at most $12$ in
$\{-1,0,1\}^{12}$. Exact enumeration of this box gives only the
origin in that closed ball. The primary certificate instead uses
rational completion of squares, visiting $59$ nodes with the same
result. The separate implementation
\path{independent/principal_unmarked_bc_v122.py} reconstructs $Q$
by cyclotomic companion matrices and checks all $3^{12}=531441$ box
vectors by an integer split quadratic form. Both paths verify the
unimodular change of basis and the complete bounds.
In fact the lattice minimum is $14$: $H$ is even integral, the
closed ball of norm $12$ contains only zero, and its first basis vector
has norm $14$.

Thus the necessary trace norm $12$ is impossible, excluding every
polarized isomorphism $B_B\simeq B_C$. This leaves their established
unmarked isogeny relation unchanged. The stronger unpolarized distinction
is proved in Proposition~\ref{pp:prop:unpolarizeddistinct} below; neither
argument asserts injectivity of the whole principal curve or identifies
it with Mostaed's named image.
\end{proof}


\subsection{Scalar-polarized maps}
\begin{proposition}[The smallest scalar-polarized $B$--$C$ isogeny]
\label{pp:prop:bcisogeny}
For the specified principal fibers $(B_B,\lambda_B)$ and
$(B_C,\lambda_C)$, the smallest positive integer $k$ for which an
isogeny $f:B_B\to B_C$ satisfies $f^*\lambda_C=k\lambda_B$ is $2$.
There are exactly six such isogenies at $k=2$, each of degree $64$.
They have the same kernel, a maximal isotropic subgroup of $B_B[2]$
isomorphic to $(\Z/2\Z)^6$, and form one orbit under the polarized
automorphisms of $B_C$. The origin-preserving polarized automorphism
groups of both fibers are the explicit cyclic groups of order six
generated, in their respective CM coordinates, by multiplication by $z^7$.
\end{proposition}
\begin{proof}
Use $I_B,I_C,\xi_B,\xi_C,r,J,Q,H$ and the unique Hodge twist
$\sigma_{31}$ from Proposition~\ref{pp:prop:unmarkeddistinct}.
A rational Hodge map $m\sigma_{31}$ is an isogeny with multiplier $k$
if and only if
\[
 m\sigma_{31}(I_B)\subseteq I_C,\qquad m\bar m=kr.
\]
The inclusion must be tested on the full integral lattices; membership
in $J$ alone is insufficient. The norm equation is nonzero, so it makes
the rational map invertible. An integral Hodge map between these
abelian varieties is algebraic, as in the homology comparison above.
On the common integral symplectic form $E$, its matrix satisfies
$T^tET=kE$, so Pfaffians give its degree $\det T=k^6$.
Moreover, any rational scalar multiplier is integral: unimodularity
makes $E$ primitive, so integer linear combinations of its entries
equal $1$, and the corresponding combinations of the entries of
$T^tET$ equal $k$. Positivity then gives $k>0$.

For $k=2$ the necessary trace norm is $24$. The exact inverse-diagonal
bounds for $H$ give
\[
 (|n_1|,\ldots,|n_{12}|)\leq(1,1,1,1,2,1,1,1,1,1,1,1).
\]
The certificate \path{certificates/principal_bc_isogeny_v123.py}
and the separate companion-matrix implementation
\path{independent/principal_bc_isogeny_v123.py} prove completeness
by rational suffix recursion and Schur-complement prefix minima,
respectively. The independent path records $175$ vectors in the closed
trace-norm ball of value $24$, and both find $84$ on its boundary.
Of the latter, $72$ fail the
full relative norm equation, and six more fail lattice integrality.
The remaining six satisfy the full norm equation, all $144$ matrix
integrality conditions, $T^tET=2E$ and $\det T=64$.

Explicitly, one map is $F_0(v)=m_0\sigma_{31}(v)$, where
\[
\begin{split}
127m_0={}&19-98z-85z^2-37z^3+15z^4+48z^5-56z^6\\
        &{}+24z^7+7z^8+105z^9+82z^{10}+2z^{11}.
\end{split}
\]
Its integral matrix is retained in
\path{certificates/principal_bc_isogeny_data.json}; the implementations
reconstruct it directly from the actual bases of $I_B,I_C$ and verify
these identities. The complete six-map list is
$F_j=z^{7j}F_0$, $0\leq j<6$. Multiplication by $z^7$ preserves each
actual integral lattice and its polarization and has exact order six;
both implementations check these statements directly.
Thus all six maps have the same kernel.

The identity $T^tET=2E$ gives the integral matrix
$2T^{-1}=E^{-1}T^tE$. Hence the cokernel of $T$ is killed by $2$;
its order $64$ makes it $(\Z/2\Z)^6$. It identifies with the kernel
of the isogeny in $B_B[2]$. For representatives $v,w$ of kernel points,
$2E_B(v,w)=E_C(F_0v,F_0w)$ is integral, so their $2$-Weil pairing is
trivial. Its order $2^6$ makes this isotropic subgroup maximal.
The independent certificate also computes the Smith form
$\operatorname{diag}(1,1,1,1,1,1,2,2,2,2,2,2)$ for every map.

Postcomposition with any polarized automorphism of $B_C$ permutes the
six maps faithfully, since $F_0$ is surjective. Its orbit of $F_0$
therefore has the size of that automorphism group and at most six
elements. The six displayed units already occur, proving that they
are the full group. The same argument by precomposition applies to
$B_B$: if $F_0a=F_0b$, the connected image of the homomorphism $a-b$
lies in the finite kernel of $F_0$, so $a=b$.

Finally multiplier one would have degree one and give a polarized
isomorphism, excluded by Proposition~\ref{pp:prop:unmarkeddistinct}.
Thus multiplier two is minimal. This does not assert a minimum degree
among arbitrary unpolarized isogenies, nor a fixed-$\OL$-linear map;
the twist $\sigma_{31}$ changes the prescribed real multiplication.
\end{proof}

\subsection{Unpolarized isomorphisms and minimum degree}
\begin{proposition}[Three distinct underlying abelian varieties]
\label{pp:prop:unpolarizeddistinct}
The three specified principal models $B_A,B_B,B_C$ of
Theorem~\ref{thm:principalfamily} are pairwise nonisomorphic as complex
abelian varieties, even after forgetting their polarizations and real
multiplication markings.
\end{proposition}
\begin{proof}
The accepted simple-factor decomposition distinguishes $B_A$ from
$B_B,B_C$. For the latter pair, use the actual CM-coordinate lattices
$I_B,I_C$ and the rational Hodge twist $\sigma_{31}$ from
Proposition~\ref{pp:prop:unmarkeddistinct}. Write
$S_X=\OM I_X$ for their maximal-order saturations, and let
$F_0=m_0\sigma_{31}$ be the explicit degree-$64$ isogeny of
Proposition~\ref{pp:prop:bcisogeny}. Put
\[
 s=z^6+z^{12}+z^{24}-z^{18}-z^{30}-z^{36},\qquad
 \alpha=\frac{1+s}{2},\qquad \varpi=1-z^{14}.
\]
Then $s^2=-7$, $\alpha\bar\alpha=2$, and exact lattice calculations give
\begin{equation}\label{pp:eq:unpolarizedlocaldata}
 F_0(S_B)=\alpha S_C,\qquad [S_X:I_X]=27,\qquad
 [\OM:R]=27,\quad R=(I_C:I_C),\quad
 v-\bar v\in\varpi\OM\quad(v\in R).
\end{equation}
The first equality includes an integral determinant-one basis change;
an equality of indices alone would not suffice. The last condition is
verified on a basis of the \emph{full} multiplier order, not merely a
selected suborder.

Here is a finite exact description of the computation. In the basis
$1,z,\ldots,z^{11}$, the saturation is the lattice generated by
$z^i I_X$, $0\leq i<12$. For each basis element of $M$, write its
multiplication matrix on $I_C$ and collect the $144$ entries as a column
of a rational matrix $A$. The full multiplier order consists exactly
of the coefficient vectors $x$ for which $Ax\in\Z^{144}$; it is the
dual of the lattice generated by the rows of $A$.
The certificate \path{certificates/principal_bc_unpolarized_v125.py}
reconstructs these lattices by integer Hermite forms after clearing
denominators. The separate certificate
\path{independent/principal_bc_unpolarized_v125.py} checks the
saturations by containment, stability under $z$ and generation modulo
$3$, and the full multiplier order by linear integrality constraints
modulo $3$. Both verify \eqref{pp:eq:unpolarizedlocaldata} on the actual
C56 lattices and the displayed $F_0$.

The ideal $\mathfrak P=\varpi\OM$ is the prime above $3$ with residue
field $\mathbb F_{729}$. Indeed $\zeta_3=z^{14}$ reduces to $1$, and
$3$ has order $6$ modulo $7$, so $\Phi_7$ is irreducible over
$\mathbb F_3$. Complex conjugation acts on this field by $x\mapsto
x^{27}$, since $27\equiv-1\pmod7$.

Suppose now that $F=m\sigma_{31}$ were an unpolarized isomorphism
$B_B\to B_C$. It would send $I_B$ onto $I_C$ and hence $S_B$ onto
$S_C$. Comparing this with \eqref{pp:eq:unpolarizedlocaldata} shows that
\[
 \eta S_C=S_C,\qquad \eta=\alpha m/m_0.
\]
Multiplication by $\eta$ and by $\eta^{-1}$ therefore has an integral
matrix on the full lattice $S_C$. Their characteristic polynomials
show that both are algebraic integers: $\eta\in\OM^\times$.
This uses no class-number-one or principal-ideal choice.

Let $\Z_{(3)}$ denote localization at the rational prime $3$.
Since $\det F_0=64$, the maps $F$ and $F_0$ are both isomorphisms
from $I_B\otimes\Z_{(3)}$ to $I_C\otimes\Z_{(3)}$.
Thus multiplication by $u=m/m_0$ belongs to
$(R\otimes\Z_{(3)})^\times$. By the last condition in
\eqref{pp:eq:unpolarizedlocaldata}, its nonzero residue is fixed by
complex conjugation. Consequently
\[
 \frac{\eta}{\bar\eta}\equiv
 \frac{\alpha}{\bar\alpha}\equiv s\pmod{\mathfrak P}.
\]
The second congruence uses $s^2=-1$ in characteristic $3$:
$(1+s)/(1-s)=s$. The residue of $s$ has exact order $4$.
On the other hand, the CM unit $\eta/\bar\eta$ has absolute value $1$
at every complex embedding, so Kronecker's theorem makes it a root
of unity in $M$, hence an element of $\mu_{42}$. In characteristic
$3$ its reduction has order dividing $14$. The two orders are
incompatible, proving the contradiction.
The degree-$64$ isogenies remain valid. The next proposition determines
the least degree among all unpolarized isogenies.
\end{proof}

\begin{proposition}[The least unpolarized isogeny degree]
\label{pp:prop:mindegree}
Among all isogenies $B_B\to B_C$ of the specified principal models,
with no polarization compatibility imposed, the least degree is $7$.
There is an explicit such isogeny with cyclic kernel of order $7$.
\end{proposition}
\begin{proof}
Retain $I_X,S_X,m_0,\alpha$ from the preceding propositions and put
$\beta=1+z-z^3$. Its absolute norm is $7$. Define
\[
 m_7=z^5\beta\alpha^{-1}m_0,\qquad F_7=m_7\sigma_{31}.
\]
Explicitly,
\[
\begin{split}
127m_7={}&-40-41z-55z^2-9z^3+62z^4+46z^5+31z^6\\
       &{}+23z^7+12z^8+53z^9+68z^{10}-51z^{11}.
\end{split}
\]
The matrix $I_C^{-1}M(m_7)\sigma_{31}(I_B)$ is integral with
determinant $7$. Thus the rational Hodge map $F_7$ induces an algebraic
isogeny of degree $7$, with cyclic kernel. Its exact integral matrix is
retained in \path{certificates/principal_bc_min_degree_data.json}.
The certificate \path{certificates/principal_bc_min_degree_v126.py}
checks every entry and the determinant; the separate companion
implementation \path{independent/principal_bc_min_degree_v126.py}
reconstructs the map independently and checks the Smith form.

For the lower bound, any isogeny $F=m\sigma_{31}$ sends $I_B$ into
$I_C$, and hence sends $S_B$ into $S_C$. The exact identity
$F_0(S_B)=\alpha S_C$ gives
\[
 F(S_B)=\eta S_C\subseteq S_C,\qquad \eta=\alpha m/m_0.
\]
Its integral multiplication matrix makes $\eta$ an algebraic integer.
Since $[S_B:I_B]=[S_C:I_C]=27$, lattice indices yield
\[
 \deg F=[I_C:F(I_B)]=[S_C:F(S_B)]
       =[S_C:\eta S_C]=N_{M/\mathbb Q}(\eta).
\]
The last norm is positive because $M$ is CM. If it were between $2$
and $6$, the ideal $(\eta)$ would have a prime ideal factor of norm
$2$, $3$, $4$ or $5$. None exists. Indeed $\OM=\Z[z]$, and exact
polynomial arithmetic gives
\[
 \gcd(\Phi_{42},X^4-X)=1\quad\text{over }\mathbb F_2,
 \qquad
 \gcd(\Phi_{42},X^p-X)=1\quad\text{over }\mathbb F_p,
 \quad p=3,5.
\]
These exclude residue degrees one or two above $2$ and degree one
above $3$ or $5$, exhausting the possible prime-power norms at most
$6$. The primary certificate checks the three gcds; the independent
certificate evaluates $\Phi_{42}$ on $\mathbb F_4$, $\mathbb F_3$
and $\mathbb F_5$ and verifies that there are no roots.
Degree one is excluded by Proposition~\ref{pp:prop:unpolarizeddistinct}.
Together with $F_7$, this proves the asserted minimum.

The map $F_7$ is not scalar-polarized for the specified principal
polarizations: a positive rational scalar multiplier would be an
integer $k$ and give degree $k^6$, which cannot equal $7$.
This is compatible with the minimum scalar multiplier $2$ and degree
$64$ in Proposition~\ref{pp:prop:bcisogeny}. We assert no finite census
of all unpolarized isogenies.
\end{proof}

\subsection{The integral Hom ideal}
\begin{proposition}[Integral maps and an ideal class of order two]
\label{pp:prop:homideal}
Set $R=(I_C:I_C)=\OL+\mathfrak P$, where
$\mathfrak P=(1-z^{14})\OM$, and define
\[
 \mathfrak a=\alpha\OL+\mathfrak P,
 \qquad \gamma=z^6(1-z^2).
\]
All integral homomorphisms $B_B\to B_C$ are, uniquely,
\[
 F_\eta=\frac{m_0}{\alpha}\eta\sigma_{31},\qquad
 \eta\in\mathfrak a.
\]
The normalized module $\mathfrak a$ is an invertible fractional
$R$-ideal under target postcomposition, with
\[
 [\OM:\mathfrak a]=27,
 \qquad \mathfrak a\bar{\mathfrak a}=R,
 \qquad \mathfrak a^2=\gamma R,
 \qquad [\mathfrak a]\in\operatorname{Pic}(R)
       \text{ of exact order }2.
\]
For every nonzero $\eta$, the map $F_\eta$ is an isogeny of degree
$N_{M/\mathbb Q}(\eta)$. If
$E_{B,q}(v,w)=\operatorname{Tr}_{M/\mathbb Q}(\xi_B qv\bar w)$,
then its actual pullback on the source lattice is
\[
 F_\eta^*E_C=E_{B,\,\sigma_{19}(\eta\bar\eta)}.
\]
In particular the degree-$7$ map of Proposition~\ref{pp:prop:mindegree}
has pullback polarization type $(1,1,1,1,1,7)$.
\end{proposition}
\begin{proof}
Every rational Hodge map has the unique form $m\sigma_{31}$.
The map $(m_0/\alpha)\sigma_{31}$ identifies $S_B$ with $S_C$
by Proposition~\ref{pp:prop:unpolarizeddistinct}.
Thus an integral map must have $\eta S_C\subseteq S_C$, so
$\eta\in\OM$. Away from $3$, both original lattices equal their
saturations locally, because their saturation indices are $27$;
there is no further local restriction there.
At $3$, the map $F_0$ is a lattice isomorphism and $\alpha$ is a unit.
The remaining condition is exactly
\[
 \eta/\alpha\in R\otimes\Z_{(3)}.
\]
Since $R=\OL+\mathfrak P$ has residue field $\mathbb F_{27}$
inside $\OM/\mathfrak P=\mathbb F_{729}$, this says that the
residue of $\eta$ lies in the line $\alpha\mathbb F_{27}$.
Its full inverse image in $\OM$ is precisely $\mathfrak a$.
This proves the description of all maps, including sufficiency.

These module equalities also have direct finite certificates. Collecting
the $144$ entries of the map for each coefficient of $\eta$ gives
rational linear integrality constraints. The primary certificate
\path{certificates/principal_bc_hom_v127.py} takes the exact dual
of their row lattice; the result is $\mathfrak a$, not just a
sublattice of candidate maps. The independent implementation
\path{independent/principal_bc_hom_v127.py} derives the same module
from linear conditions modulo $3$, with containment and index checks.
The primary additionally checks $R=\OL+\mathfrak P$; both verify
the full products
$\mathfrak a\bar{\mathfrak a}=R$ and
$\mathfrak a^2=\gamma R$ on integral bases.
They check that $\gamma\in\OM^\times$ has absolute norm $1$.
The first product gives an explicit inverse fractional ideal.
If $\mathfrak a=\delta R$ were principal, the equality
$[\OM:\mathfrak a]=[\OM:R]=27$ would give
$N_{M/\mathbb Q}(\delta)=1$. Since $\delta\in\mathfrak a$,
the corresponding map would be an isomorphism, contradicting
Proposition~\ref{pp:prop:unpolarizeddistinct}. Hence the class is
nontrivial, and the second product proves that its order is exactly two.
The full Picard group is determined in Proposition~\ref{pp:prop:picardorder} below.

The degree formula is the saturation-index argument of
Proposition~\ref{pp:prop:mindegree}. For the polarization, the identity
$m_0\bar m_0=2r$ and $\alpha\bar\alpha=2$ gives
$m\bar m=r\eta\bar\eta$. Moving $\sigma_{31}$ through the trace
therefore gives
\[
 E_C(F_\eta v,F_\eta w)
 =\operatorname{Tr}_{M/\mathbb Q}
   \bigl(\xi_B\sigma_{19}(\eta\bar\eta)v\bar w\bigr),
\]
since $\sigma_{19}=\sigma_{31}^{-1}$. The inverse twist is essential
on source coordinates. For nonzero $\eta$ its multiplier is totally
positive; integrality follows from the integral map.
The independent checker verifies all $78$ quadratic coefficients
of this identity on a basis of $\mathfrak a$, including the mixed terms.
For the degree-$7$ witness, both paths check the alternating Smith
invariants $(1^{10},7^2)$, giving the displayed polarization type.
The formula describes exactly the pullbacks arising from this Hom
module; it does not assert that every totally positive element is realized.
\end{proof}

\subsection{Ideal classes and principal polarizations}
\begin{proposition}[The full Picard group of the principal-model order]
\label{pp:prop:picardorder}
For the order $R=\OL+\mathfrak P$ in
Proposition~\ref{pp:prop:homideal},
\[
 \operatorname{Pic}(R)\cong\mathbb Z/2\mathbb Z,
 \qquad
 \operatorname{Pic}(R)=\{[R],[\mathfrak a]\},
 \qquad \mathfrak a=\alpha\OL+\mathfrak P.
\]
Here $\operatorname{Pic}(R)$ denotes invertible fractional $R$-ideals
modulo principal fractional ideals, without an orientation or positivity
condition.
\end{proposition}
\begin{proof}
Write $O=\OM$. The conductor and residue computations above give
\[
 O/\mathfrak P=\mathbb F_{729},\qquad
 R/\mathfrak P=\mathbb F_{27}.
\]
Masley's Main Theorem gives class number one for
$M=\mathbb Q(\zeta_{42})=\mathbb Q(\zeta_{21})$:
$\varphi(21)=12\le24$, and $21$ is absent from its exceptional table
\cite[p.~179]{Masley1976}. This concerns the maximal order $O$;
the following argument determines
the class group of the proper order $R$.

Let $I$ be an invertible fractional $R$-ideal. Since $IO$ is principal,
rescale $I$ so that $IO=O$. Then
\[
 \mathfrak P\subseteq I\subseteq O:
 \qquad \mathfrak P=\mathfrak P(IO)\subseteq I
\]
because $\mathfrak P O\subseteq R$.
Away from the conductor maximal ideal $\mathfrak P$ of $R$,
the order is maximal and the normalized ideal equals $O$ locally.
At $\mathfrak P$, invertibility gives
$I_{\mathfrak P}=uR_{\mathfrak P}$, where normalization makes
$u$ a unit of $O_{\mathfrak P}$. Therefore $I/\mathfrak P$ is a
nonzero $\mathbb F_{27}$-line in $\mathbb F_{729}$.
Conversely, the inverse image $J_a$ of any such line
$a\mathbb F_{27}$ is locally principal: at $\mathfrak P$ it is
$uR_{\mathfrak P}$ for a lift $u$ of $a$, and elsewhere it is $O$.
Thus every such $J_a$ is invertible, and
$J_aJ_b=J_{ab}$ by the same local calculation.

Any isomorphism between fractional ideals is multiplication by an
element of $M^\times$. Between normalized ideals it must preserve
their $O$-saturations, so it is multiplication by a global $O$-unit.
Consequently this classifies every invertible ideal class and gives
\[
 \operatorname{Pic}(R)\cong
 \frac{\mathbb F_{729}^{\times}}
 {\mathbb F_{27}^{\times}\,\operatorname{im}(O^{\times})}.
\]
The intermediate quotient
$Q=\mathbb F_{729}^{\times}/\mathbb F_{27}^{\times}$ is cyclic
of order $28$. Conjugation induces $x\mapsto x^{27}$ on the residue
field, and the ratio map
\[
 Q\longrightarrow\mathbb F_{729}^{\times},
 \qquad [x]\longmapsto x/x^{27}
\]
is injective, since its kernel consists of the fixed-field units.

For every $\varepsilon\in O^{\times}$, Kronecker's theorem gives
$\varepsilon/\bar\varepsilon\in\mu_{42}$. The reduction of
$\mu_{42}$ has order $14$, so the image of the entire global unit
group in $Q$ is contained in the subgroup of order $14$.
It equals that subgroup. Indeed $z$ has residue order $14$,
and its ratio $z^2$ has order $7$. The global unit
$\gamma=z^6(1-z^2)$ has norm $1$ and
\[
 \gamma/\bar\gamma=-z^{14}\equiv-1\pmod{\mathfrak P},
\]
so its class has order $2$. These two classes generate a subgroup
of order $14$. The remaining quotient therefore has order $2$.

Finally, $\alpha\bar\alpha=2$ makes the residue of $\alpha$
nonzero. With $s=\sqrt{-7}$ chosen as above,
\[
 \alpha/\bar\alpha\equiv s\pmod{\mathfrak P},
 \qquad s^2\equiv-1\pmod{\mathfrak P}.
\]
Its ratio has order $4$, so its class lies outside the global-unit
image of order $14$. Its residue line has inverse image
$\mathfrak a=\alpha\OL+\mathfrak P$. This identifies the
nonidentity class and agrees with the ideal product identities of
Proposition~\ref{pp:prop:homideal}.

The finite inputs are checked by
\path{certificates/principal_order_class_v129.py}, using Frobenius
and polynomial gcds over $\mathbb F_3$ and exact rational resultants.
The separate implementation
\path{independent/principal_order_class_v129.py} checks all $39$
possible monic divisors of degrees at most three, enumerates the
$729$ residue elements, and computes the global norms by rational
multiplication matrices. Both identify the residue field using
\[
 \Phi_{42}(X)\bmod3
 =(X^6-X^5+X^4-X^3+X^2-X+1)^2.
\]
These finite checks supplement the class-number theorem and the
ideal and global-unit arguments above; they do not enumerate
the infinite group of global units.
\end{proof}

\begin{proposition}[The ordinary ideal-class monoid of the selected order]
\label{pp:prop:idealclassmonoid}
Let $O=\OM$ and retain $R$ and $\mathfrak a$ from
Proposition~\ref{pp:prop:picardorder}. Let $\operatorname{ICM}(R)$
be the monoid of nonzero fractional $R$-ideals that are full
finitely generated $\mathbb Z$-lattices in $M$, modulo $R$-linear
isomorphism, with ordinary ideal multiplication. Then
\[
 \operatorname{ICM}(R)=\{[R],[\mathfrak a],[O]\},\qquad
 [\mathfrak a]^2=[R],\quad [\mathfrak a][O]=[O]^2=[O].
\]
Its unit group is $\operatorname{Pic}(R)\cong\mathbb Z/2\mathbb Z$;
$[O]$ is the remaining, absorbing idempotent. The multiplier rings
$(I:I)=\{x\in M:xI\subseteq I\}$ are $R$ for the first two
classes and $O$ for the third. These are ordinary rank-one classes:
the zero ideal is excluded, and equivalence is $R$-linear, without
orientation or semilinear identifications. This statement for the
selected order makes no geometric realization or polarization claim.
\end{proposition}
\begin{proof}
Use the conductor $\mathfrak P$, residue fields
$k=R/\mathfrak P=\mathbb F_{27}$ and
$k_2=O/\mathfrak P=\mathbb F_{729}$, and
$\operatorname{Pic}(O)=1$ established in the preceding proposition.
For any ideal $I$ in the stated class, $IO$ is a nonzero fractional
ideal of the Dedekind domain $O$. It is principal, even when $I$
is not invertible over $R$. A scalar rescaling therefore gives
$IO=O$ and $I\subseteq O$. Write $1=\sum_j i_j o_j$ with
$i_j\in I$ and $o_j\in O$. For $p\in\mathfrak P$,
\[
 p=\sum_j i_j(po_j)\in I,
\]
since $po_j\in\mathfrak P\subseteq R$. Thus
$\mathfrak P\subseteq I\subseteq O$, and $I/\mathfrak P$
is a nonzero $k$-subspace of $k_2$; otherwise $I=\mathfrak P$
would contradict $IO=O$. Since $[k_2:k]=2$, this subspace is
either a line or all of $k_2$. The latter case gives $I=O$.

For $t\in k_2^\times$, let $J_t$ be the inverse image of $tk$.
A lift of $t$, together with $\mathfrak P$, generates $O$ as an
$O$-ideal, so $J_tO=O$. The product $J_tJ_u$ also has saturation
$O$ and hence contains $\mathfrak P$ by the same argument.
Its residue is exactly $tuk$, giving
\[
 J_tJ_u=J_{tu},\qquad J_1=R,\qquad J_tJ_{t^{-1}}=R.
\]
Every line lift is therefore invertible over $R$.

An $R$-linear isomorphism between these full rank-one ideals
extends over $\operatorname{Frac}(R)=M$ to multiplication by
$c\in M^\times$. Between normalized ideals it satisfies $cO=O$,
so $c\in O^\times$. The preceding proposition gives exactly two
line classes, represented by $R$ and $\mathfrak a$. The class of
$O$ is distinct, since an $O$-unit cannot carry $O$ to a proper
line lift. Normalization gives $[I][O]=[O]$ for every class,
and the remaining products follow from $\operatorname{Pic}(R)$.

Finally, $xI\subseteq I$ for a normalized ideal implies
$xO\subseteq O$, hence $x\in O$. For $I=J_t$, preservation of
the residue line $tk$ is equivalent to $x\bmod\mathfrak P\in k$,
so $(J_t:J_t)=R$; also $(O:O)=O$. These rings are unchanged by
scalar rescaling. As $R\ne O$, the third class is not invertible:
an inverse ideal $H$ would give $xR=xIH\subseteq IH=R$ for
every $x\in(I:I)$, forcing $(I:I)=R$.
\end{proof}

\begin{proposition}[Marked principal polarizations for the selected order]
\label{pp:prop:markedpolarizations}
Put $O=\OM$, and retain $R$ and $\mathfrak a$ from
Proposition~\ref{pp:prop:homideal}. Fix a marked CM type $\Phi$ of $M$,
obtained from $\Phi_0=\{1,5,11,13,17,19\}$ by conjugating a subset
$T$ of its six embeddings. Among full fractional $R$-ideals, up to
$M$-linear polarized isomorphism, the principal polarizations whose
Rosati involution on $M$ is complex conjugation are as follows:
\[
\begin{array}{c|ccc|c}
 & [R]&[\mathfrak a]&[O]&\text{total at }\Phi\\ \hline
|T|\text{ even}&2&2&0&4\\
|T|\text{ odd}&0&0&1&1
\end{array}
\]
The origin-preserving marked polarized automorphism groups are
$\mu_6$ in the first row and $\mu_{42}$ in the second. Thus there
are $160$ marked polarized classes over the $64$ marked CM types.
This count makes no identification between different types and no
assertion that these objects lie on a specified curve.
\end{proposition}
\begin{proof}
Write $I^\vee=\{x\in M:\Tr_{M/\Ql}(xI)\subseteq\Z\}$.
An $M$-compatible alternating form is
$E_\xi(x,y)=\Tr_{M/\Ql}(\xi x\bar y)$, with $\bar\xi=-\xi$.
Its dual lattice is
\[
 I^{\#_\xi}=\xi^{-1}(\bar I)^\vee
             =\xi^{-1}(R^\vee:\bar I).
\]
Indeed $x\bar I\subseteq R^\vee$ is equivalent to
$\Tr(x\bar I)\subseteq\Z$, since $\bar I R=\bar I$.
For invertible $I$, this colon equals $R^\vee\bar I^{-1}$.
Replacing $R^\vee$ by the maximal-order codifferent here would be
incorrect for a nonmaximal order.

Set $d=f'(t)$, retain $\beta$ from Proposition~\ref{pp:prop:explicit},
and put
\[
 \kappa_O=\frac1{d(z-\bar z)},\qquad
 \kappa_R=\frac1{d\beta(z-\bar z)}.
\]
The identity $\beta O=(1-z^{14})O$ gives
$O=\OL\oplus\OL z$ and $R=\OL\oplus\OL\beta z$.
Monogenic relative trace duality, followed by trace transitivity,
therefore gives $O^\vee=\kappa_O O$ and $R^\vee=\kappa_R R$.
Since $\mathfrak a\bar{\mathfrak a}=R$, one also has
$(\bar{\mathfrak a})^\vee=\kappa_R\mathfrak a$.
Consequently principal compatible parameters on $R,\mathfrak a$
are $\kappa_R u$, and those on $O$ are $\kappa_O u$, with
$u\in E_L=\OL^\times$. For the nonmaximal models this follows
because the real units of their multiplier ring $R$ are exactly $E_L$.
The conductor identity is certified by
\path{certificates/cm_polarization_conductor_v135.py}: it verifies
$\beta=(1-z^{14})a_0$ with
$a_0=1+z+z^5+z^6-z^8+z^{11}$ an integral unit of norm one.

With the convention $E(x,Jx)>0$, positivity means
$\operatorname{Im}\phi(\xi)>0$ for every $\phi\in\Phi$.
At $\Phi_0$, the signs of $d$ are $(+,-,+,-,+,-)$ and
$\operatorname{Im}\phi(z)>0$. Thus $\kappa_O$ has odd negative
sign parity. Since $N_{L/\Ql}(\beta)=-27$, $\kappa_R$ has even
parity. Conjugating an embedding reverses its imaginary sign.
Lemma~\externalref{supplement}{lem:narrowclass} gives exactly the even-parity subgroup
as the real-unit signature image, and
$[E_L^+:E_L^2]=2$. This proves the asserted existence and
nonexistence of positive parameters.

We next compute the relevant unit norms. For $e\in R^\times$,
Kronecker's theorem gives $e/\bar e\in\mu_{42}$.
Its residue modulo $\mathfrak P$ is one because $R/\mathfrak P$
is the conjugation-fixed field $\mathbb F_{27}$.
The reduction of $\mu_{42}$ has order $14$, so its kernel is
$\mu_3$. Choose $\zeta\in\mu_6$ with $\zeta^2=e/\bar e$.
Then $e/\zeta\in E_L$. Conversely $\mu_6E_L\subset R^\times$,
giving
\[
 R^\times=\mu_6E_L,\qquad N_{M/L}(R^\times)=E_L^2.
\]
By Theorem~\externalref{supplement}{thm:unitnorm}, $N_{M/L}(O^\times)=E_L^+$.
For a fixed positive parameter the polarization classes form a torsor
under $E_L^+/N_{M/L}(S^\times)$, where $S=(I:I)$.
Proposition~\ref{pp:prop:idealclassmonoid} exhausts the three scalar
classes and their multiplier rings, so these quotients prove the table.
Polarized automorphisms have norm one and are roots of unity.
The roots in $R$ are exactly $\mu_6$, while those in $O$ are
$\mu_{42}$. Finally $32\cdot4+32\cdot1=160$.
\end{proof}

\begin{proposition}[All principal polarizations on the specified B/C models]
\label{pp:prop:actualpolarizations}
Let $B_B,B_C$ be the two specified principal models, with their
original polarizations $\lambda_B,\lambda_C$. Put
$\varepsilon_+=4-t^2$. Their scalar ideal classes are respectively
$[\mathfrak a]$ and $[R]$. On each underlying variety $B_X$,
exactly two isomorphism classes of principal polarizations occur,
represented by $\lambda_X$ and $\lambda_X\circ\iota(\varepsilon_+)$.
After forgetting all CM and RM markings these still give exactly
four polarized isomorphism classes over the two underlying varieties.
Each has origin-preserving polarized automorphism group $\mu_6$.
\end{proposition}
\begin{proof}
Use the exact $I_B,I_C,\xi_B,\xi_C$ coordinates of
Proposition~\ref{pp:prop:unmarkeddistinct}, and $m_0,\alpha$ from
Proposition~\ref{pp:prop:homideal}. The scalars
\[
\begin{split}
 c_B={}&8+11z+4z^2+4z^3+8z^4-5z^5-22z^6-24z^7\\
       &{}-2z^8+20z^9+18z^{10}+3z^{11},\qquad
 c_C=\alpha\sigma_{31}(c_B)/m_0
\end{split}
\]
satisfy $c_BI_B=\mathfrak a$ and $c_CI_C=R$.
The certificate \path{certificates/actual_bc_polarizations_v135.py}
verifies integral unimodular changes of basis for both equalities,
and transports each form to
$\xi'_X=\xi_X/(c_X\bar c_X)=\kappa_Ru_X$ with $u_X\in E_L$.
It checks the exact alternating matrices and their pullbacks against
the original principal matrices. The complete coefficients and basis
changes are retained in
\path{certificates/actual_bc_polarizations_data.json}.
These are finite identities using one arithmetic implementation;
the following classification is a human argument.

There is also a structural reason for $c_C$.
For $D=(m_0/\alpha)\sigma_{31}(I_B)$, the complete Hom calculation
gives $(I_C:D)=\mathfrak a$. The proved B identification makes
$D=(m_0/\alpha)\sigma_{31}(c_B)^{-1}\bar{\mathfrak a}$ invertible;
here $\sigma_{31}(\mathfrak a)=\bar{\mathfrak a}$.
Multiplying by $D$ and using
$\mathfrak a\bar{\mathfrak a}=R$ gives the C identification.

For $\gamma=z^6(1-z^2)$ one has
\[
 \varepsilon_+=\gamma\bar\gamma=4-t^2,\qquad
 \gamma/\bar\gamma=-z^{14}\notin\mu_{21}.
\]
Thus $\varepsilon_+$ is a totally positive real unit.
If it were $v^2$ with $v\in E_L$, Kronecker's theorem would give
$\gamma/v\in\mu_{42}$, forcing
$\gamma/\bar\gamma\in\mu_{42}^2=\mu_{21}$, a contradiction.
It represents the nontrivial class of $E_L^+/E_L^2$.
Proposition~\ref{pp:prop:markedpolarizations} now supplies the two
compatible principal polarization classes on each model.

These are all principal polarizations, even without a marking.
The actual B and C types in Proposition~\ref{pp:prop:markeddistinct}
each have trivial stabilizer in $\operatorname{Gal}(M/\Ql)$.
The disjoint permutation supports of
$\operatorname{End}_{\Ql}(M)=\bigoplus_\sigma M\sigma$
then show $\operatorname{End}^0(B_X)=M$.
For any polarization $\lambda=\lambda_Xa$, its Rosati involution is
$x\mapsto a^{-1}\bar x a=\bar x$, since $M$ is commutative.
Thus every polarization is compatible, and every automorphism lies
in the multiplier group $R^\times$. Its norm subgroup is $E_L^2$,
so the two classes remain distinct without markings.
Proposition~\ref{pp:prop:unpolarizeddistinct} excludes all cross
identifications between B and C. The automorphism groups follow
from Proposition~\ref{pp:prop:markedpolarizations}.
No corresponding claim is made for the induced A type.
\end{proof}

